\section{Les oscillateurs}

Un oscillateur est un montage électronique délivrant un signal périodique.La fréquence peut généralement être réglée via le choix des composants.

Deux grandes familles d'oscillateurs existent :
\begin{itemize}
  \item les oscillateurs quasi-sinusoïdaux produisent des signaux proches de sinusoïdes
  \item les oscillateurs à relaxation produisent des signaux créneaux ou triangulaires
\end{itemize}

Les oscillateurs sont souvent des systèmes bouclés.

\section{Oscillateur quasi-sinusoïdal : l'oscillateur de Wien}
\subsection{Montage de l'oscillateur de Wien}

Les oscillateurs quasi-sinusoïdaux sont très souvent constitués
\begin{itemize}
  \item d'un amplificateur
  \item d'un filtre passe-bande
\end{itemize}

Pour le filtre de Wien, le filtre passe-bande est un filtre de Wien et l'amplificateur est un amplificateur non-inverseur.

\formule{Montage de l'oscillateur de Wien}[]{\schema{}{5cm}}[%
\begin{itemize}
  \item $v_1(t)$ est la sortie de l'amplificateur non-inverseur et l'entrée du filtre de Wien (en \unit{V})
  \item $v_2(t)$ est l'entrée de l'amplificateur non-inverseur et la sortie du filtre de Wien (en \unit{V})
\end{itemize}][o]

\flashcard{Montage de l'oscillateur de Wien.}{}

\formule{Schéma bloc}[%
  \begin{itemize}
    \item la circuit est dans l'ARQS
    \item la fréquence est dans la bande passante de l'amplificateur non-inverseur
    \item l'ALI est en régime linéaire
  \end{itemize}]{\schema{}{4cm}}[%
  \begin{itemize}
    \item $R_1$, $R_2$ et $R$ sont des résistances (en \unit{\ohm})
    \item $C$ est la capacité des condensateurs (en \unit{F})
    \item $V_1(p)$ est la sortie de l'amplificateur non-inverseur et l'entrée du filtre de Wien (en \unit{V})
    \item $V_2(p)$ est l'entrée de l'amplificateur non-inverseur et la sortie du filtre de Wien (en \unit{V})
  \end{itemize}
  ][][o]

\colle{Citer le montage et établir le schéma-bloc de l'oscillateur de Wien.}

\subsection{Condition d'oscillation et fréquence des oscillations}

\formule{Condition d'oscillation}[%
  \begin{itemize}
    \item la circuit est dans l'ARQS
    \item la fréquence est dans la bande passante de l'amplificateur non-inverseur
    \item l'ALI est en régime linéaire
  \end{itemize}]{Des oscillations sinusoïdales existent ssi $R_2=2R_1$}[%
  \begin{itemize}
    \item $R_1$ et $R_2$ sont les résistances de l'amplificateur non-inverseur (en \unit{\ohm})
  \end{itemize}][][o]

\formule{Pulsation des oscillations sinusoïdales}[%
  \begin{itemize}
    \item la circuit est dans l'ARQS
    \item la fréquence est dans la bande passante de l'amplificateur non-inverseur
    \item l'ALI est en régime linéaire
    \item la condition d'oscillation est vérifiée
  \end{itemize}]{$$\omega=\frac{1}{RC}$$}[%
  \begin{itemize}
    \item $R$ est la résistance des résistors dans le filtre de Wien (en \unit{\ohm})
    \item $C$ est la capacité des condensateurs dans le filtre de Wien (en \unit{F})
  \end{itemize}][][o]

De manière générale, un oscillateur quasi-sinusoïdal composé d'un filtre passe-bande et d'un amplificateur oscille à la pulsation caractéristique du filtre passe-bande si les conditions d'oscillation sont vérifiées.

\colle{Le schéma-bloc de l’oscillateur de Wien étant fourni, établir la condition d’existence d’oscillations sinusoïdales ainsi que leur fréquence.}

\subsection{Démarrage des oscillations}

Si le système est stable et que les tensions sont initialement nulles, rien ne se passe. Pour que des oscillations démarrent, le système bouclé doit être instable\footnote{C'est le système bouclé qui doit être instable. Chacun des deux systèmes qui le composent sont stables.}.

\formule{Condition de démarrage des oscillations}[
  \begin{itemize}
    \item la circuit est dans l'ARQS
    \item la fréquence est dans la bande passante de l'amplificateur non-inverseur
    \item l'ALI est en régime linéaire
  \end{itemize}
]
{Les oscillations démarrent ssi $R_2>2R_1$}
[
  \begin{itemize}
    \item $R_1$ et $R_2$ sont les résistances de l'amplificateur non-inverseur (en \unit{\ohm})
  \end{itemize}
][][o]

La condition d'existence d'oscillations sinusoïdales apparait comme un cas limite de cette inégalité.

\subsection{Saturation de l'ALI}

S'il n'y avait pas la saturation de l'ALI, l'amplitude des oscillations continuerait à augmenter indéfiniment. C'est la saturation de l'ALI qui fixe l'amplitude des oscillations.

\formule{Amplitude des tensions}[
  \begin{itemize}
    \item la circuit est dans l'ARQS
    \item la fréquence est dans la bande passante de l'amplificateur non-inverseur
    \item l'ALI est en régime linéaire
    \item la condition d'oscillation est vérifiée
  \end{itemize}
]
{$v_1$ a pour amplitude $V_{sat}$\\$v_2$ a pour amplitude $\frac{V_{sat}}{3}$}
[
  \begin{itemize}
    \item $v_1(t)$ la sortie de l'amplificateur non inverseur et l'entrée du filtre de Wien (en \unit{V})
    \item $v_1(t)$ l'entrée de l'amplificateur non inverseur et la sortie du filtre de Wien (en \unit{V})
    \item $V_\text{sat}$ la tension de saturation de l'ALI ($\approx\SI{15}{V}$) (en \unit{V})
  \end{itemize}
][][o]

La saturation de l'ALI est un phénomène non linéaire, qui modifie donc les spectres de $v_1$ et $v_2$. Expérimentalement, on observe que plus $R_2-2R_1$ est grand, plus les spectres comportent d'harmoniques et moins les signaux sont sinusoïdaux.

$v_2$ est <<plus sinusoïdal>> que $v_1$ car c'est la sortie du filtre passe-bande, qui diminue l'amplitude relative des harmoniques.

\colle{Le schéma-bloc de l’oscillateur de Wien étant fourni, établir la condition de démarrage des oscillations et établir l’amplitude des oscillations sinusoïdales pour les deux tensions.}

\section{Oscillateur à relaxation}
\subsection{Montage de l'oscillateur à relaxation}
L'oscillateur à relaxation est constitué d'un comparateur à hystérésis positif et d'un intégrateur.

\formule{Montage de l'oscillateur à relaxation}[]{\schema{}{5cm}}
[
  \begin{itemize}
    \item $u(t)$ est la sortie du comparateur à hystérésis positif et l'entrée de l'intégrateur (en \unit{V})
    \item $v(t)$ est l'entrée du comparateur à hystérésis positif et la sortie de l'intégrateur (en \unit{V})
  \end{itemize}
]

\flashcard{Montage de l'oscillateur à relaxation.}{}

\formule{Caractéristique du comparateur à hystérésis positif}[]{\schema{}{5cm}}
[
  \begin{itemize}
    \item $u(t)$ est la sortie du comparateur à hystérésis positif (en \unit{V})
    \item $v(t)$ est l'entrée du comparateur à hystérésis positif (en \unit{V})
  \end{itemize}
][][o]

\formule{Fonction intégrateur}
[
  \begin{itemize}
    \item le circuit est dans l'ARQS
    \item l'ALI est en régime linéaire
  \end{itemize}
]
{
  $$v(t)=v_0-\frac{1}{RC}\int_0^t u(x)dx$$
}
[
  \begin{itemize}
    \item $v(t)$ est la sortie de l'intégrateur (en \unit{V})
    \item $u(t)$ est l'entrée de l'intégrateur (en \unit{V})
    \item $R$ est la résistance du résistor (en \unit{\ohm})
    \item $C$ est la capacité du condensateur (en \unit{F})
    \item $v_0=v(t=0)$ est la valeur initiale de $v$ (en \unit{V})
  \end{itemize}
][][o]

\colle{Citer le montage de l’oscillateur à relaxation et établir l'équation différentielle du montage intégrateur et la caractéristique du montage comparateur à hystérésis positif.}

\subsection{Signaux de sortie}

\formule{Alure des signaux de sortie}
[
  \begin{itemize}
    \item le circuit est dans l'ARQS
    \item l'ALI de l'intégrateur est en régime linéaire
    \item la période est très grande devant la durée de commutation de l'ALI
  \end{itemize}
]
{\schema{}{4cm}}
[
  \begin{itemize}
    \item $V_{sat}$ est la tension de saturation des ALI (en \unit{V})
    \item $R_1$ et $R_2$ les résistances des résistors du comparateur à hystérésis positif (en \unit{\ohm})
    \item $u(t)$ est la sortie du comparateur à hystérésis positif et l'entrée de l'intégrateur (en \unit{V})
    \item $v(t)$ est l'entrée du comparateur à hystérésis positif et la sortie de l'intégrateur (en \unit{V})
  \end{itemize}
][o][o]

\subsection{Période d'oscillation}

\formule{Période d'oscillation de l'oscillateur à relaxaton}
[
  \begin{itemize}
    \item le circuit est dans l'ARQS
    \item l'ALI de l'intégrateur est en régime linéaire
    \item la période est très grande devant la durée de commutation de l'ALI
  \end{itemize}
]
{$$T=4\frac{R_1}{R_2}RC$$}
[
  \begin{itemize}
    \item $R_1$ et $R_2$ les résistances des résistors du comparateur à hystérésis positif (en \unit{\ohm})
    \item $R$ la résistance du résistor de l'intégrateur (en \unit{\ohm})
    \item $C$ la capacité du condensateur de l'intégrateur (en \unit{F})
  \end{itemize}
][o][o]

Cette expression est valable tant que la période est très grande devant la durée de commutation de l'ALI.
\schema{Influence de la vitesse de balayage sur le signal créneau.}{4cm}

\application{Déterminer la durée de commuation de l'ALI.}

\colle{L'équation différentielle de l'intégrateur et la caractéristique du comparateur à hystérésis étant données, établir la forme des signaux dans un oscillateur à relaxation et leur période.}

\subsection{Choix de $R_1$ et $R_2$}

$v(t)$ étant la sortie de l'ALI, $v\in [-V_{sat},V_{sat}]$. Pour que la commutation ait lieu, il faut que $\frac{R_1}{R_2}V_{sat}<V_{sat}$, i.e. que $R_1<R_2$.